% 声明为子文件，指定主文件
\documentclass[main.tex]{subfiles}

\begin{document}
\pagestyle{plain}
\setcounter{chapter}{0}

\chapter{Directional Patterns}
\label{chap:chapter1}

After stabilizing 3000 trials of $301\times 301$ sandpile grids (each with background 3, plus 3 extra grains at random sites), we examine the pattern formed by cells whose height differs from 3. Far from the center, along a fixed rational direction $(n,m)$, a periodic height pattern emerges. We classify these patterns for all directions with $0 \le n,m \le 7$.

\section{Classification Table for $(n,m)$, $0 \le n,m \le 7$}

Table~\ref{tab:direction-grid} shows the directional pattern for each pair $(n,m)$. Only seven distinct primitive (coprime) directions were observed across all 3000 trials. When a pair $(n,m)$ reduces by $\gcd(n,m)$ to one of these seven primitive directions, the corresponding image is shown; otherwise the entry reads ``not found.''

\begin{table}[H]
\centering
\caption{Directional pattern classification for $0 \le n,m \le 7$. Rows index $n$, columns index $m$. Entries show the image of the reduced coprime direction.}
\label{tab:direction-grid}
\small
\setlength{\tabcolsep}{3pt}
\begin{tabular}{|c|c|c|c|c|c|c|c|c|}
\hline
$n \downarrow \, / \, m \to$ & $0$ & $1$ & $2$ & $3$ & $4$ & $5$ & $6$ & $7$ \\
\hline
$0$ & --- & \includegraphics[width=1.1cm]{0_1} & \includegraphics[width=1.1cm]{0_1} & \includegraphics[width=1.1cm]{0_1} & \includegraphics[width=1.1cm]{0_1} & \includegraphics[width=1.1cm]{0_1} & \includegraphics[width=1.1cm]{0_1} & \includegraphics[width=1.1cm]{0_1} \\
\hline
$1$ & \includegraphics[width=1.1cm]{1_0} & \includegraphics[width=1.1cm]{1_1} & \includegraphics[width=1.1cm]{1_2} & \includegraphics[width=1.1cm]{1_3} & not found & not found & not found & not found \\
\hline
$2$ & \includegraphics[width=1.1cm]{1_0} & \includegraphics[width=1.1cm]{2_1} & \includegraphics[width=1.1cm]{1_1} & not found & \includegraphics[width=1.1cm]{1_2} & not found & \includegraphics[width=1.1cm]{1_3} & not found \\
\hline
$3$ & \includegraphics[width=1.1cm]{1_0} & \includegraphics[width=1.1cm]{3_1} & not found & \includegraphics[width=1.1cm]{1_1} & not found & not found & \includegraphics[width=1.1cm]{1_2} & not found \\
\hline
$4$ & \includegraphics[width=1.1cm]{1_0} & not found & \includegraphics[width=1.1cm]{2_1} & not found & \includegraphics[width=1.1cm]{1_1} & not found & not found & not found \\
\hline
$5$ & \includegraphics[width=1.1cm]{1_0} & not found & not found & not found & not found & \includegraphics[width=1.1cm]{1_1} & not found & not found \\
\hline
$6$ & \includegraphics[width=1.1cm]{1_0} & not found & \includegraphics[width=1.1cm]{3_1} & \includegraphics[width=1.1cm]{2_1} & not found & not found & \includegraphics[width=1.1cm]{1_1} & not found \\
\hline
$7$ & \includegraphics[width=1.1cm]{1_0} & not found & not found & not found & not found & not found & not found & \includegraphics[width=1.1cm]{1_1} \\
\hline
\end{tabular}
\end{table}

The seven observed directions are:
\[
(0,1),\;(1,0),\;(1,1),\;(1,2),\;(1,3),\;(2,1),\;(3,1).
\]
Note that $(n,m)$ and $(m,n)$ yield distinct patterns in general (e.g., $(1,2)$ vs.\ $(2,1)$ are different), and the negative directions $(-n,m)$, $(n,-m)$, $(-n,-m)$ produce the same patterns by symmetry.

\section{Observed Regularities}

Beyond the classification, several quantitative regularities were noted across all trials.

\subsection{Internal Color Periodicity}

Along any directional line, the color pattern of heights repeats periodically without variation. For instance, a line with  $(2, 1)$ always follows the sequence (brown, white, brown) from left to right.

\subsection{Frequency of Occurrence}

Lines whose slope is close to horizontal or vertical appear less frequently across trials. The vast majority of observed lines are horizontal, vertical, or at a $45^\circ$ angle (slope~$\pm1$). Lines with slopes $1/2$, $1/3$, $2$, or~$3$ are comparatively rare.

\subsection{Vector Sum at Intersections}

When multiple directional lines meet at a single intersection point, the sum of their unit direction vectors equals zero. This suggests a form of local conservation governing the pattern geometry.


\end{document}